\documentclass[11pt]{article}

\usepackage[margin=1in]{geometry}
\usepackage{amsmath,amssymb,mathtools,bm}
\usepackage{booktabs}
\usepackage{graphicx}
\usepackage{microtype}
\usepackage{hyperref}
\usepackage[nameinlink,capitalise]{cleveref}
\usepackage[numbers,sort&compress]{natbib}

\newcommand{\E}{\mathbb{E}}
\newcommand{\R}{\mathbb{R}}
\newcommand{\ind}{\mathbb{I}}
\newcommand{\CS}{\mathrm{CS}}

\title{Counterfactual Susceptibility:\\Discovering Inactive and Inhibitory Circuits in Language Models}
\author{Anonymous Authors}
\date{}

\begin{document}
\maketitle

\begin{abstract}
Sparse attribution graphs typically describe features that are active on a baseline prompt. This leaves out counterfactual mechanisms in which an inactive feature is held below threshold by an active inhibitory pathway but becomes causally important after a small intervention. We introduce \emph{counterfactual susceptibility}, a prompt-local score that combines a target feature's activation margin with the response of its preactivation to a specified upstream intervention. We distinguish ease of activation from predicted behavioral salience and propose intervention sweeps and mediation tests as ground truth. This manuscript is a research scaffold: no empirical claims are made until the method is validated in both replacement and underlying models.
\end{abstract}

\section{Introduction}
\label{sec:introduction}

Mechanistic interpretability aims to identify internal variables and causal pathways that explain model behavior. Sparse feature dictionaries and attribution graphs make this problem more tractable by replacing dense neural activations with a smaller set of active, potentially interpretable features \citep{ameisen2025circuit,lindsey2025biology}. Their sparsity is also a structural limitation: the graph for a prompt is normally built from features active on that prompt.

An inactive feature need not be causally irrelevant. It may lie just below its activation threshold because one or more active upstream features inhibit it. Suppressing an upstream feature can move the inactive feature across its gate, recruit a previously invisible pathway, and change the model's output. An active-only graph cannot display this pathway before the intervention.

We study the following prediction problem: given a baseline prompt, an active source feature, an inactive target feature, and a declared intervention family, can we predict (i) whether the target will activate, (ii) the intervention strength at which it activates, and (iii) whether the new activation mediates a behavior change?

Our proposed framework separates three claims that are often conflated. \emph{Activation susceptibility} measures how close a feature is to being recruited by a particular intervention. \emph{Behavioral salience} estimates the consequence of that recruitment for a target metric. \emph{Causal mediation} tests whether the recruited feature explains part of the source intervention's behavioral effect.

The intended contributions are:
\begin{enumerate}
    \item a dimensionless, prompt-local susceptibility score derived from activation margins and source-to-target responses;
    \item a candidate-generation procedure for searching large inactive dictionaries without exhaustive intervention;
    \item a benchmark based on continuous intervention sweeps, measured critical intervention strengths, and held-out prompts;
    \item mediation experiments that distinguish behaviorally meaningful counterfactual circuits from incidental gate crossings;
    \item an augmentation of active-only attribution graphs with empirically verified inactive nodes and inhibitory edges.
\end{enumerate}

This version defines the hypotheses and experimental protocol. Results will be added only after the full pipeline is reproduced and audited.

\section{Background}
\label{sec:background}

\subsection{Sparse replacement models and attribution graphs}

Cross-layer transcoders replace dense MLP computations with sparse feature activations whose decoder outputs can write to later layers. A local replacement model can then be used to construct prompt-specific feature interaction graphs \citep{ameisen2025circuit}. These graphs are useful hypothesis generators, but their relation to the underlying model must be checked by interventions.

\subsection{Thresholded features}

Let feature $i$ have preactivation $z_i$ and threshold $\tau_i$, with activation
\begin{equation}
    a_i=\phi_i(z_i),
\end{equation}
where $\phi_i$ is the exact nonlinearity used by the loaded dictionary. We deliberately leave the implementation-specific JumpReLU convention abstract here; the empirical code must read thresholds, biases, and activation semantics from the pinned model implementation.

\subsection{The active-only blind spot}

If $a_i=0$, feature $i$ is normally absent from the baseline attribution graph. Yet a perturbation can alter $z_i$, change the active set, and recruit feature $i$ into the downstream computation. This is especially important for inhibitory mechanisms and for evaluating local linear approximations under interventions.

\section{Counterfactual Susceptibility}
\label{sec:method}

\subsection{Single-source suppression}

Fix a prompt and an inactive target feature $i$ with positive margin
\begin{equation}
    m_i=\tau_i-z_i>0.
\end{equation}
Let $j$ be an active upstream feature with activation $a_j$. We consider suppressing its activation by a fraction $\alpha\in[0,1]$.
Our intervention convention is
\begin{equation}
    a_j \mapsto (1-\alpha)a_j,
\end{equation}
so $\alpha=0$ leaves the source unchanged and $\alpha=1$ fully suppresses it.

Let
\begin{equation}
    J_{ij}=\frac{\partial z_i}{\partial a_j}
\end{equation}
be the prompt-local response of the target preactivation to the source activation under a declared linearization convention. The first-order counterfactual target preactivation is
\begin{equation}
    \widehat z_i(\alpha)=z_i-\alpha a_jJ_{ij}.
\end{equation}

Define the movement toward the target threshold under full source suppression as
\begin{equation}
    q_{j\rightarrow i}=-a_jJ_{ij}.
\end{equation}
For an active source with positive activation, an inhibitory local relation has
$J_{ij}<0$, hence $q_{j\rightarrow i}>0$: suppressing the source raises the
predicted target preactivation. Conversely, $q_{j\rightarrow i}\leq 0$ does not
move the target toward its gate under this intervention. These statements are
prompt-local, first-order predictions; finite interventions may change active
sets and other nonlinear model states, so crossings must be verified by an
intervention sweep.
For $q_{j\rightarrow i}>0$, the predicted critical intervention is
\begin{equation}
    \widehat\alpha^\star_{j\rightarrow i}
    =\frac{m_i}{q_{j\rightarrow i}}.
\end{equation}
A full suppression reaches or exceeds the threshold in the local-linear
preactivation when $\widehat\alpha^\star_{j\rightarrow i}\leq 1$. At equality,
the prediction lands exactly at $z_i=\tau_i$; whether this activates the feature
depends on whether the loaded nonlinearity uses a strict or non-strict threshold
comparison. We therefore treat equality, and values within numerical tolerance
of it, as boundary cases rather than definite crossings.

\subsection{Dimensionless susceptibility}

We define pairwise counterfactual susceptibility as
\begin{equation}
    S_{j\rightarrow i}
    =\frac{q_{j\rightarrow i}}{m_i+\varepsilon}.
    \label{eq:susceptibility}
\end{equation}
For the unregularized score ($\varepsilon=0$), and away from numerical tolerance
boundaries, $S_{j\rightarrow i}>1$ is equivalent to
$\widehat\alpha^\star_{j\rightarrow i}<1$. With $\varepsilon>0$, however,
$S_{j\rightarrow i}>1$ requires $q_{j\rightarrow i}>m_i+\varepsilon$ and is
therefore a conservative sufficient condition for a strict full-suppression
crossing, not an exact equivalence.

The predicted target activation at intervention strength $\alpha$ is
\begin{equation}
    \widehat a_i(\alpha)
    =\phi_i\!\left(z_i-\alpha a_jJ_{ij}\right).
\end{equation}

\subsection{Behavior-weighted salience}

Let $T$ be a scalar target metric and
\begin{equation}
    g_i=\frac{\partial T}{\partial a_i}.
\end{equation}
We define a first-order behavior-weighted score
\begin{equation}
    B_{j\rightarrow i}(\alpha)
    =|g_i|\,|\widehat a_i(\alpha)-a_i|.
\end{equation}
We evaluate $S$ and $B$ separately: activation ease is not behavioral importance.

\subsection{Causal mediation}

For a verified source--target pair, let $\Delta T_j$ denote the behavior change under source suppression and let $\Delta T_{j,i0}$ denote the change when the source is suppressed while the target is clamped inactive. The signed mediated effect is
\begin{equation}
    \Delta T^{\mathrm{med}}_{j\rightarrow i}
    =\Delta T_j-\Delta T_{j,i0}.
\end{equation}
When numerically stable, we also report the mediated fraction
\begin{equation}
    \mathrm{MF}_{j\rightarrow i}
    =\frac{\Delta T^{\mathrm{med}}_{j\rightarrow i}}{\Delta T_j}.
\end{equation}

\subsection{Candidate generation}

A full dictionary search is computationally prohibitive. We therefore use a staged search: retain near-threshold inactive features, filter for positive predicted movement under selected source suppressions, and optionally prioritize features with large downstream behavioral response. Candidate recall is audited against a smaller brute-force intervention subset.

\section{Experimental Design}
\label{sec:experiments}

\subsection{Evaluation ladder}

We proceed from synthetic verification to open-model experiments. First, a small thresholded feature network with known inhibitory edges tests the formulas and numerical edge cases. Second, an open language model with pretrained transcoders is used to measure the prevalence of intervention-induced gate crossings. Third, held-out prompts are used to compare susceptibility against preregistered baselines. Finally, high-ranked pairs are tested with behavioral and mediation interventions.

\subsection{Ground-truth intervention sweeps}

For each selected active source, we sweep $\alpha$ from zero to one, record target preactivations and activations, bracket any gate crossing, and estimate the observed critical $\alpha^\star$ by binary search. We record replacement-model and underlying-model consequences separately.

\subsection{Baselines}

The primary baselines are random inactive targets, margin alone, inhibitory response alone, downstream gradient alone, and an unthresholded product of response and gradient. On a smaller subset, exhaustive intervention provides an oracle discovery reference.

\subsection{Metrics}

Gate-crossing prediction is measured with AUPRC, precision at fixed candidate budgets, and candidate-generation recall. Critical-strength prediction is measured with rank correlation, absolute error, and calibration. Behavioral predictions are evaluated by sign accuracy and normalized error. Mediation experiments report signed mediated effects and matched random-target controls. Efficiency is reported in interventions, accelerator time, and peak memory.

\subsection{Generalization and controls}

Discovery and evaluation prompts will be disjoint. Any key results will be checked across random seeds, dictionary sizes, and, if feasible, a second model family. Planned controls include matched random subspaces or features, non-inhibitory source--target pairs, and intervention strengths below the predicted threshold.

\section{Limitations}
\label{sec:limitations}

The score is prompt-local and initially first order. Interventions can change feature gates, attention patterns, normalization terms, and downstream compensatory mechanisms, invalidating the local response. Sparse dictionary features may split or absorb concepts, and replacement-model responses may differ from those of the underlying model. Candidate filtering can miss important far-from-threshold features activated by nonlinear or multi-source effects. We therefore treat susceptibility as a ranking hypothesis, not a mechanistic conclusion, until verified by intervention.

\section{Related Work}
\label{sec:related}

This project builds on sparse autoencoder and transcoder methods, cross-layer replacement models, attribution graphs, activation patching, causal tracing, and feature intervention. Anthropic's circuit-tracing work explicitly identifies inactive, suppressed features as a limitation of active-only attribution graphs and suggests searching for features that are one ablation away from activation \citep{ameisen2025circuit}.

Counterfactual interventions can reveal causal dependencies but can also produce ambiguous effects or miss relevant causes when the intervention is poorly chosen \citep{mueller2024counterfactuals}. We therefore use continuous source-suppression sweeps, report minimum verified intervention strengths, and separate predictions in the local replacement model from observations in the underlying model.

Self-repair and nonlinear interactions can cause both gradient and ablation methods to underestimate component importance \citep{edin2025gim}. These effects motivate explicit calibration of local susceptibility predictions against intervention ground truth rather than treating a first-order score as a conclusion.

Benchmarks with known circuits provide a complementary way to test whether a discovery method retrieves true causal structure \citep{gupta2024interpbench}. Our planned initial synthetic tests will serve this role for threshold-crossing and inhibitory-edge recovery; the main empirical target remains sparse features in open language models.

The final review will distinguish counterfactual inactive-feature discovery from active-node attribution, contrastive activation analysis, feature steering, sensitivity analysis, and global circuit localization.

\section{Conclusion}
\label{sec:conclusion}

We propose counterfactual susceptibility as an operational framework for identifying inactive sparse features that may become causally relevant after inhibitory pathways are perturbed. The decisive evidence will be prospective prediction of gate crossings, calibrated critical intervention strengths, and mediation of behavior in the underlying model. This manuscript will be updated as those experiments are completed.


\bibliographystyle{plainnat}
\bibliography{bibliography}

\appendix
\section{Reproducibility Checklist}
The final version will report code revisions, model and dictionary revisions, configurations, random seeds, hardware, prompt-generation procedures, and complete metric tables.

\end{document}
